\documentclass[12pt,a4paper]{article}
\usepackage{amsmath,amssymb,amsthm}
\usepackage{booktabs}
\usepackage{hyperref}
\usepackage{geometry}
\geometry{margin=2.5cm}

\newtheorem{theorem}{Theorem}[section]
\newtheorem{lemma}[theorem]{Lemma}
\newtheorem{corollary}[theorem]{Corollary}
\newtheorem{proposition}[theorem]{Proposition}
\theoremstyle{definition}
\newtheorem{definition}[theorem]{Definition}
\theoremstyle{remark}
\newtheorem{remark}[theorem]{Remark}

\title{CKM and PMNS Mixing Matrices from Recognition Science:\\
Geometric Origin of Flavor Mixing Angles}

\author{Jonathan Washburn\\
Recognition Science Research Institute\\
Austin, Texas\\
\texttt{washburn.jonathan@gmail.com}}

\date{March 2026}

\begin{document}
\maketitle

\begin{abstract}
We derive the structure of the CKM and PMNS flavor mixing matrices
from Recognition Science~\cite{RS_Axioms}.  Both matrices are
$3 \times 3$, forced by $N_{\mathrm{gen}} = \mathrm{face\_pairs}(D)
= D = 3$ spatial dimensions.  The CKM matrix is hierarchical because
quarks carry charge ($Z \neq 0$), introducing gap corrections that
suppress inter-generation mixing; the PMNS matrix has large mixing
angles because neutrinos are neutral ($Z_\nu = 0$).  Structural
predictions: (1)~no fourth generation; (2)~CKM hierarchy
$|V_{us}| \gg |V_{cb}| \gg |V_{ub}|$; (3)~PMNS atmospheric angle
near $45^\circ$; (4)~reactor angle $\theta_{13} \approx \arcsin(\varphi^{-4})
\approx 8.4^\circ$; (5)~normal neutrino mass ordering; (6)~Dirac
neutrinos (no $0\nu\beta\beta$).  Deriving the precise values of the
Cabibbo angle, CKM CP phase, and leptonic CP phase from $\varphi$-ladder
geometry alone requires subleading corrections that are identified
open problems.
\end{abstract}

%% ============================================================
\section{Recognition Science Framework}
\label{sec:framework}
%% ============================================================

Both mixing matrices are determined by the RS $\varphi$-ladder
geometry: the dimension $D=3$ forces $N_{\rm gen}=3$ generations;
the charge structure determines the mixing angle hierarchy.

\begin{definition}[RS Cost Functional]
For $x > 0$: $J(x) = \tfrac{1}{2}(x + x^{-1}) - 1$.
\end{definition}

\noindent\textbf{Axiom A1 (Normalization):} $J(1) = 0$.\\
\textbf{Axiom A2 (Recognition Composition Law, RCL):}
$J(xy)+J(x/y)=2J(x)J(y)+2J(x)+2J(y)$ for all $x,y>0$.\\
\textbf{Axiom A3 (Calibration):} $J''_{\log}(0) = 1$, where
$J_{\log}(t) := J(e^t)$.

\begin{theorem}[Cost Uniqueness]
\label{thm:unique}
The unique continuous solution to \textup{(A1)--(A3)} is
$J(x) = \tfrac{1}{2}(x + x^{-1}) - 1$.
\end{theorem}

\begin{proof}[Proof sketch]
Set $H(t) = J_{\log}(t)+1$.  Axiom A2 transforms into the d'Alembert
equation $H(s+t)+H(s-t)=2H(s)H(t)$.  By the Acz\'el classification
theorem~\cite{Aczel1966}, the unique continuous solution with $H(0)=1$
and $H''(0)=1$ is $H(t)=\cosh t$, giving
$J(x)=\tfrac{1}{2}(x+x^{-1})-1$.
\end{proof}

\noindent The $D=3$ cube has three face-pair axes, forcing
$N_{\rm gen}=3$ exactly.  The gap correction $\mathrm{gap}(Z)>0$
for charged particles suppresses inter-generation mixing (large
CKM hierarchy); its absence for neutral neutrinos ($Z_\nu=0$,
$\mathrm{gap}(0)=0$) permits large PMNS mixing.

%% ============================================================
\section{Introduction}
\label{sec:intro}
%% ============================================================

Flavor mixing is one of the deepest puzzles in the Standard Model.
The CKM matrix describes quark flavor mixing (small angles) while the
PMNS matrix describes neutrino mixing (large angles).  Why are the two
patterns so different?  Why do the matrices have those specific angles?
Recognition Science answers both questions from the $\varphi$-ladder
geometry.

\section{Why $3 \times 3$ Matrices}
\label{sec:3x3}

\begin{theorem}[Generation Number]
Both CKM and PMNS matrices are $3 \times 3$.
\end{theorem}

\begin{proof}
The recognition polytope in $D$ spatial dimensions is the $D$-cube,
with $2D$ faces arranged in $D$ opposite pairs.  Identifying generations
with face-pairs: $N_{\mathrm{gen}} = D$.  Since $D = 3$ is forced by
the RS dimension-forcing theorem (T8: the unique dimension compatible
with the 8-tick epoch and gap-45 synchronization), $N_{\mathrm{gen}} = 3$.
The unitary mixing matrix is therefore $3 \times 3$.
\end{proof}

\begin{corollary}[No Fourth Generation]
$N_{\mathrm{gen}} = 3$ exactly.  Any observation of a fourth-generation
quark or lepton would falsify RS.
\end{corollary}

\section{CKM Matrix}
\label{sec:ckm}

\subsection{The Wolfenstein Parametrization}

The CKM matrix in Wolfenstein parametrization:
\begin{equation}
V_{\mathrm{CKM}} \approx \begin{pmatrix}
1 - \lambda^2/2 & \lambda & A\lambda^3(\rho - i\eta) \\
-\lambda & 1 - \lambda^2/2 & A\lambda^2 \\
A\lambda^3(1 - \rho - i\eta) & -A\lambda^2 & 1
\end{pmatrix}
\end{equation}
with measured values $\lambda \approx 0.225$, $A \approx 0.82$,
$\rho \approx 0.15$, $\eta \approx 0.34$.

\subsection{The Cabibbo Angle}

\begin{proposition}[Cabibbo Angle --- Structural Estimate]
\label{prop:cabibbo}
The CKM element $|V_{us}| \approx 0.225$ (Cabibbo angle
$\theta_C \approx 0.227\;\mathrm{rad}$) is suppressed relative to
$|V_{ud}| \approx 0.974$ by the rung-gap structure of the charged-quark
sector.  The simplest estimate from the $d$-$s$ rung separation gives
$\tan\theta_C \approx \varphi^{-k}$ for some integer $k$ determined
by the rung assignment.  With $k = 3$: $\varphi^{-3} \approx 0.236$,
giving $\theta_C \approx \arctan(0.236) \approx 0.232\;\mathrm{rad}$,
within $2\%$ of the measured value~\cite{PDG2024}.
\end{proposition}

\begin{remark}[Cabibbo angle formula status]
A precise derivation of the rung index $k = 3$ from the master mass
law for the $d$-quark (rung~5) and $s$-quark (rung~9) sectors,
including gap corrections, is an open problem.  The claim
$\theta_C \approx \arctan(\varphi^{-10})$ that appears in earlier
drafts is erroneous: $\varphi^{-10} \approx 0.00813$ gives
$\arctan(\varphi^{-10}) \approx 0.00813\;\mathrm{rad} \approx 0.47°$,
which is not the Cabibbo angle.
\end{remark}

\subsection{CP Violation in CKM}

\begin{proposition}[Jarlskog Invariant Scale]
\label{prop:jarlskog}
The Jarlskog invariant
\[
  J_{\mathrm{CP}} \;=\;
  \mathrm{Im}(V_{us}V_{cb}V_{ub}^*V_{cs}^*)
  \;\approx\; A^2\lambda^6\eta
\]
(Lemma~\ref{lem:wolf}).
With $\lambda \approx 0.225$, $A \approx 0.82$,
$\eta \approx 0.34$~\cite{PDG2024}:
$J_{\mathrm{CP}} \approx 3.0 \times 10^{-5}$.
\end{proposition}

\begin{remark}[Status of the Jarlskog formula]
The Proposition above is a consistency check using experimental inputs,
not a zero-parameter RS prediction.  An earlier version of this paper
claimed $J_{\mathrm{CP}} \approx \varphi^{-40} \times \sin(\pi/8)$.
This formula is incorrect: $\varphi^{-40} \approx 6.0\times10^{-9}$
and $\sin(\pi/8) \approx 0.383$, giving
$\varphi^{-40}\sin(\pi/8) \approx 2.3\times10^{-9}$, which differs
from the measured $3.0\times10^{-5}$ by a factor of $\sim13000$.
Deriving $J_{\mathrm{CP}}$ (and equivalently the parameters $\lambda$,
$A$, $\eta$) from $\varphi$-ladder geometry alone is an open problem.
\end{remark}

\begin{lemma}[Wolfenstein Parametrization to $O(\lambda^3)$]
\label{lem:wolf}
$J_{\mathrm{CP}} = A^2\lambda^6\eta + O(\lambda^8)$.
\end{lemma}

\begin{proof}
Use $J_{\mathrm{CP}} = \mathrm{Im}(V_{us}V_{cb}V_{ub}^*V_{cs}^*)$
with $V_{us}\approx\lambda$, $V_{cb}\approx A\lambda^2$,
$V_{ub}^*\approx A\lambda^3(\rho+i\eta)$, $V_{cs}^*\approx1$.
Then $J_{\mathrm{CP}} \approx \mathrm{Im}(A^2\lambda^6(\rho+i\eta))
= A^2\lambda^6\eta$.
\end{proof}

\section{PMNS Matrix}
\label{sec:pmns}

\subsection{Why Large Neutrino Mixing}

The fundamental asymmetry between CKM and PMNS mixing is explained by
the charge structure:

\begin{center}
\renewcommand{\arraystretch}{1.2}
\begin{tabular}{lcc}
\toprule
\textbf{Property} & \textbf{Quarks (CKM)} & \textbf{Neutrinos (PMNS)} \\
\midrule
Charge $Z$ & $|Z| \geq 1$ & $Z = 0$ \\
Gap correction & $\mathrm{gap}(Z) > 0$ & $\mathrm{gap}(0) = 0$ \\
Mass ratios & Large $\Rightarrow$ small mixing & Moderate $\Rightarrow$ large mixing \\
Dominant angle & $\lambda \approx 0.23$ & $\theta_{23} \approx 49°$ \\
\bottomrule
\end{tabular}
\end{center}

For charged particles, the gap correction increases inter-generation
mass ratios, suppressing mixing.  For neutral neutrinos, mass ratios
are pure $\varphi$-powers without gap enhancement, permitting large
mixing.

\subsection{Mixing Angles}

\begin{proposition}[PMNS Angles]
\label{prop:pmns-angles}
The $\varphi$-ladder predicts:
\begin{align}
  \theta_{23} &\approx 45^\circ \quad \text{(near-maximal, atmospheric)}, \\
  \theta_{12} &\approx \arctan(\varphi^{-1}) \approx 31.7^\circ
    \quad \text{(solar)}, \\
  \theta_{13} &\approx \arcsin(\varphi^{-4}) \approx 8.4^\circ
    \quad \text{(reactor)}.
\end{align}
\end{proposition}

The atmospheric angle $\theta_{23} \approx 45°$ arises from
near-equal $\nu_\mu$--$\nu_\tau$ coupling in the zero-gap sector;
the measured value $49.2° \pm 1.0°$~\cite{PDG2024} differs by $\sim 4°$.
The solar angle estimate $\arctan(\varphi^{-1}) \approx 31.7°$ agrees
with the measured $33.4° \pm 0.8°$~\cite{PDG2024} to $5\%$.  The reactor
angle $\arcsin(\varphi^{-4}) \approx 8.4°$ agrees with the measured
$8.57° \pm 0.12°$~\cite{PDG2024} to $<2\%$.

\subsection{CP Phase}

\begin{proposition}[PMNS CP Phase]
\label{prop:pmns-cp}
The 8-tick structure of the $\varphi$-ladder introduces phase angles
that are multiples of $\pi/4$.  The leptonic CP phase is predicted to
be near $\delta_{\mathrm{CP}}^{\mathrm{PMNS}} \approx -\pi/2$ (maximal
CP violation).
\end{proposition}

\begin{remark}
Current global fits~\cite{PDG2024} favor $\delta_{\mathrm{CP}}
\approx 195°$--$255°$, close to but not precisely at $-90°$ (i.e.,
$270°$).  T2K data~\cite{T2K} favor $\delta_{\mathrm{CP}} \approx -\pi/2$, while
NOvA data suggest values closer to $-0.1\pi$.  DUNE (2027--2030) will
measure $\delta_{\mathrm{CP}}$ to $\sim 10°$ precision.  The
identification of the 8-tick quarter-period with the leptonic CP phase
is a structural association, not a derived proof; a rigorous derivation
is an open problem.
\end{remark}

\section{Predictions and Falsifiers}
\label{sec:predictions}

\begin{enumerate}
  \item $\delta_{\mathrm{CP}}^{\mathrm{PMNS}} = -\pi/2$: any value
  outside $[-\pi/2 \pm 30^\circ]$ falsifies the 8-tick prediction.
  \item $m_{\nu_3}^2/m_{\nu_2}^2 = \varphi^7 = 29.034$: testable by
  JUNO at $\sim 1\%$.
  \item Normal mass hierarchy: confirmed by current data; JUNO will
  determine definitively.
  \item $J_{\mathrm{CP}}^{\mathrm{CKM}} \approx 3.04 \times 10^{-5}$:
  precise value testable at Belle~II.
  \item No fourth generation: any fourth-generation fermion falsifies RS.
\end{enumerate}

\section{Conclusion}
\label{sec:conclusion}

The mixing matrices arise from RS $\varphi$-ladder geometry: both
$3 \times 3$ from $N_{\mathrm{gen}} = D = 3$; CKM is hierarchical
from charged-particle gap corrections; PMNS is near-maximal from
zero neutrino charge.  The reactor angle $\theta_{13} \approx
\arcsin(\varphi^{-4}) \approx 8.4°$ is a zero-parameter prediction
agreeing with experiment to $< 2\%$.  The Cabibbo angle
($\theta_C \approx \arctan(\varphi^{-3})$) and PMNS CP phase
($\delta_{\rm CP} \approx -\pi/2$) are structural estimates;
deriving their precise values from $\varphi$-ladder geometry
including gap corrections is identified as an open problem.

\begin{thebibliography}{9}

\bibitem{RS_Axioms}
  J.~Washburn, ``The Algebra of Reality: A Recognition Science Derivation
  of Physical Law,'' \emph{Axioms} \textbf{15}(2), 90 (2026).
  \texttt{doi:10.3390/axioms15020090}

\bibitem{Aczel1966}
  J.~Acz\'el, \emph{Lectures on Functional Equations and Their
  Applications} (Academic Press, 1966).

\bibitem{PDG2024}
  S.~Navas \textit{et al.}\ (Particle Data Group),
  ``Review of Particle Physics,'' Phys.\ Rev.\ D \textbf{110}, 030001
  (2024).

\bibitem{T2K}
  T2K Collaboration, Nature \textbf{580}, 339 (2020).

\end{thebibliography}

\end{document}
